% --- Archon physics formalization source begin ---
% archon:physics
% archon:covers PhyXMiniProblems/problem_phyx_mini_0733.lean
% archon:source-report reports/phyx_mini/problem_phyx_mini_0733.source.json

\chapter{Physics problem phyx\_mini\_0733}
\label{ch:PhyXMiniProblems_problem_phyx_mini_0733}

\paragraph{Problem source.}
\#\# Physical scenario

In figure, a heavy piece of machinery is raised by sliding it a distance \$d = 12.5\textbackslash{} m\$ along a plank oriented at angle \$\textbackslash{}theta = 20.0\textasciicircum{}\textbackslash{}circ\$ to the horizontal.

\#\# Question

How far is it moved vertically?

\#\# Answer choices

- A: A: 2.42 m

- B: B: 3.81 m

- C: C: 4.28 m

- D: D: 5.63 m

\#\# Figure caption (auxiliary; use the image as primary evidence)

The image depicts a schematic of an inclined plane scenario. Here's a detailed breakdown:

1. **Objects**:
   - A ramp or inclined plane tilted against the ground.
   - A cylindrical object, possibly a wheel or spool, resting on top of the inclined plane.

2. **Angles and Dimensions**:
   - The angle between the inclined plane and the horizontal ground is labeled as \textbackslash{}(\textbackslash{}theta\textbackslash{}).
   - The distance along the inclined plane from the cylindrical object to the bottom of the ramp is marked as \textbackslash{}(d\textbackslash{}).

3. **Relationships**:
   - The inclined plane is supported at an angle above the ground, and the cylindrical object is positioned near the top of the incline.
   - There is a vertical support under the inclined plane, keeping it at the set angle \textbackslash{}(\textbackslash{}theta\textbackslash{}).

4. **Surfaces**:
   - The ground is depicted with a textured pattern, likely to illustrate a rough surface.
   - The ramp is shown with a smooth surface, indicated by the clean lines representing its surface.

Overall, this image is likely illustrating a physics problem setup involving an inclined plane and the forces acting on an object on the incline.

\#\# Dataset metadata

Kinematics; Spatial Relation Reasoning

\paragraph{Recorded answer/context.}
C: C: 4.28 m

\paragraph{Figure/image path.}
/root/proposal\_for\_physic/science-mango/phyx\_mini\_run/phyx\_data/test\_image/733.png

\paragraph{Formalization target.}
create a compiling Lean file with sorry bodies at `PhyXMiniProblems/problem\_phyx\_mini\_0733.lean`. The Lean declarations must preserve the physical quantities, dimensions or dimensional roles, figure labels, governing-law hypotheses, and final relation expressed by this problem.
Use Mathlib/Physlib names found through LeanExplore where available. If a physics API is missing, introduce faithful local abstractions rather than scalar placeholder aliases.

\begin{theorem}[Physics formalization target]
\label{thm:physics:phyx_mini_0733:target}
\lean{PhyXMiniProblems.ProblemPhyXMini0733.problem_phyx_mini_0733}
\uses{def:physics:phyx-mini-0733:phyxminiproblems-problemphyxmini0733-lengthinmeters, def:physics:phyx-mini-0733:phyxminiproblems-problemphyxmini0733-anglefromdegrees, def:physics:phyx-mini-0733:phyxminiproblems-problemphyxmini0733-machineryinclinesetup, def:physics:phyx-mini-0733:phyxminiproblems-problemphyxmini0733-matchesmachineryinclinescenario, def:physics:phyx-mini-0733:phyxminiproblems-problemphyxmini0733-matchesmachineryinclinereadouts, def:physics:phyx-mini-0733:phyxminiproblems-problemphyxmini0733-matchessuppliedinclinedplankfigure, def:physics:phyx-mini-0733:phyxminiproblems-problemphyxmini0733-hasphysicalmachineryinclineparameters, def:physics:phyx-mini-0733:phyxminiproblems-problemphyxmini0733-satisfiesinclinedplankprojectionlaw, def:physics:phyx-mini-0733:phyxminiproblems-problemphyxmini0733-isuniqueroundedanswer}
The assigned autoformalize agent should translate this physics problem into Lean declarations in the covered file, with theorem and lemma proof bodies written as `by sorry`.
\end{theorem}
\begin{proof}
This is an autoformalization task, not a proof task. Produce statements that can later be proved by the physics prover without weakening the physical model.
\end{proof}

% --- Archon declaration topology begin ---
\section{Declaration topology}
\begin{definition}[Length Quantity]
  \label{def:physics:phyx-mini-0733:phyxminiproblems-problemphyxmini0733-lengthquantity}
  \lean{PhyXMiniProblems.ProblemPhyXMini0733.LengthQuantity}
  A nonnegative physical distance, represented coherently in every unit
\end{definition}
\begin{proof}
  This is the declared model definition of Length Quantity.
\end{proof}
\begin{definition}[length Readout]
  \label{def:physics:phyx-mini-0733:phyxminiproblems-problemphyxmini0733-lengthreadout}
  \lean{PhyXMiniProblems.ProblemPhyXMini0733.lengthReadout}
  \uses{def:physics:phyx-mini-0733:phyxminiproblems-problemphyxmini0733-lengthquantity}
  Read a physical distance in a selected length unit
\end{definition}
\begin{proof}
  This is the declared model definition of length Readout.
\end{proof}
\begin{definition}[length In Meters]
  \label{def:physics:phyx-mini-0733:phyxminiproblems-problemphyxmini0733-lengthinmeters}
  \lean{PhyXMiniProblems.ProblemPhyXMini0733.lengthInMeters}
  \uses{def:physics:phyx-mini-0733:phyxminiproblems-problemphyxmini0733-lengthquantity, def:physics:phyx-mini-0733:phyxminiproblems-problemphyxmini0733-lengthreadout}
  Metre readout of a physical distance
\end{definition}
\begin{proof}
  This is the declared model definition of length In Meters.
\end{proof}
\begin{definition}[angle From Degrees]
  \label{def:physics:phyx-mini-0733:phyxminiproblems-problemphyxmini0733-anglefromdegrees}
  \lean{PhyXMiniProblems.ProblemPhyXMini0733.angleFromDegrees}
  Convert a numerical degree readout to Mathlib's angle type
\end{definition}
\begin{proof}
  This is the declared model definition of angle From Degrees.
\end{proof}
\begin{definition}[Moving Body Kind]
  \label{def:physics:phyx-mini-0733:phyxminiproblems-problemphyxmini0733-movingbodykind}
  \lean{PhyXMiniProblems.ProblemPhyXMini0733.MovingBodyKind}
  Kind of body being translated along the support surface
\end{definition}
\begin{proof}
  This is the declared model definition of Moving Body Kind.
\end{proof}
\begin{definition}[Support Surface Kind]
  \label{def:physics:phyx-mini-0733:phyxminiproblems-problemphyxmini0733-supportsurfacekind}
  \lean{PhyXMiniProblems.ProblemPhyXMini0733.SupportSurfaceKind}
  Kind of surface that constrains the body's straight-line motion
\end{definition}
\begin{proof}
  This is the declared model definition of Support Surface Kind.
\end{proof}
\begin{definition}[Along Plank Direction]
  \label{def:physics:phyx-mini-0733:phyxminiproblems-problemphyxmini0733-alongplankdirection}
  \lean{PhyXMiniProblems.ProblemPhyXMini0733.AlongPlankDirection}
  Direction of translation along the plank
\end{definition}
\begin{proof}
  This is the declared model definition of Along Plank Direction.
\end{proof}
\begin{definition}[Figure Object]
  \label{def:physics:phyx-mini-0733:phyxminiproblems-problemphyxmini0733-figureobject}
  \lean{PhyXMiniProblems.ProblemPhyXMini0733.FigureObject}
  Physical objects visible in the supplied raster
\end{definition}
\begin{proof}
  This is the declared model definition of Figure Object.
\end{proof}
\begin{definition}[Figure Label]
  \label{def:physics:phyx-mini-0733:phyxminiproblems-problemphyxmini0733-figurelabel}
  \lean{PhyXMiniProblems.ProblemPhyXMini0733.FigureLabel}
  The two symbolic annotations visible in the supplied raster
\end{definition}
\begin{proof}
  This is the declared model definition of Figure Label.
\end{proof}
\begin{definition}[Supplied Inclined Plank Figure]
  \label{def:physics:phyx-mini-0733:phyxminiproblems-problemphyxmini0733-suppliedinclinedplankfigure}
  \lean{PhyXMiniProblems.ProblemPhyXMini0733.SuppliedInclinedPlankFigure}
  \uses{def:physics:phyx-mini-0733:phyxminiproblems-problemphyxmini0733-lengthquantity, def:physics:phyx-mini-0733:phyxminiproblems-problemphyxmini0733-figureobject, def:physics:phyx-mini-0733:phyxminiproblems-problemphyxmini0733-figurelabel}
  Data transcribed from image 733. The bitmap supplies incidences, orientations, and the meanings of d and theta; the numerical values come from the prose. It does not display a numerical vertical displacement
\end{definition}
\begin{proof}
  This is the declared model definition of Supplied Inclined Plank Figure.
\end{proof}
\begin{definition}[Machinery Incline Setup]
  \label{def:physics:phyx-mini-0733:phyxminiproblems-problemphyxmini0733-machineryinclinesetup}
  \lean{PhyXMiniProblems.ProblemPhyXMini0733.MachineryInclineSetup}
  \uses{def:physics:phyx-mini-0733:phyxminiproblems-problemphyxmini0733-lengthquantity, def:physics:phyx-mini-0733:phyxminiproblems-problemphyxmini0733-movingbodykind, def:physics:phyx-mini-0733:phyxminiproblems-problemphyxmini0733-supportsurfacekind, def:physics:phyx-mini-0733:phyxminiproblems-problemphyxmini0733-alongplankdirection, def:physics:phyx-mini-0733:phyxminiproblems-problemphyxmini0733-suppliedinclinedplankfigure}
  Independent quantities of the motion. In particular, verticalDisplacement is not defined from the requested answer; it is related to the path displacement by the projection law below
\end{definition}
\begin{proof}
  This is the declared model definition of Machinery Incline Setup.
\end{proof}
\begin{definition}[Matches Machinery Incline Scenario]
  \label{def:physics:phyx-mini-0733:phyxminiproblems-problemphyxmini0733-matchesmachineryinclinescenario}
  \lean{PhyXMiniProblems.ProblemPhyXMini0733.MatchesMachineryInclineScenario}
  \uses{def:physics:phyx-mini-0733:phyxminiproblems-problemphyxmini0733-machineryinclinesetup}
  Qualitative assignments stated by the problem scenario
\end{definition}
\begin{proof}
  This is the declared model definition of Matches Machinery Incline Scenario.
\end{proof}
\begin{definition}[Matches Machinery Incline Readouts]
  \label{def:physics:phyx-mini-0733:phyxminiproblems-problemphyxmini0733-matchesmachineryinclinereadouts}
  \lean{PhyXMiniProblems.ProblemPhyXMini0733.MatchesMachineryInclineReadouts}
  \uses{def:physics:phyx-mini-0733:phyxminiproblems-problemphyxmini0733-lengthinmeters, def:physics:phyx-mini-0733:phyxminiproblems-problemphyxmini0733-anglefromdegrees, def:physics:phyx-mini-0733:phyxminiproblems-problemphyxmini0733-machineryinclinesetup}
  The two numerical inputs from the prose. No vertical distance or answer choice is included here
\end{definition}
\begin{proof}
  This is the declared model definition of Matches Machinery Incline Readouts.
\end{proof}
\begin{definition}[Matches Supplied Inclined Plank Figure]
  \label{def:physics:phyx-mini-0733:phyxminiproblems-problemphyxmini0733-matchessuppliedinclinedplankfigure}
  \lean{PhyXMiniProblems.ProblemPhyXMini0733.MatchesSuppliedInclinedPlankFigure}
  \uses{def:physics:phyx-mini-0733:phyxminiproblems-problemphyxmini0733-machineryinclinesetup}
  Objects, labels, and relations read from the primary image. The label fields denote setup quantities without assigning the unknown vertical displacement
\end{definition}
\begin{proof}
  This is the declared model definition of Matches Supplied Inclined Plank Figure.
\end{proof}
\begin{definition}[Has Physical Machinery Incline Parameters]
  \label{def:physics:phyx-mini-0733:phyxminiproblems-problemphyxmini0733-hasphysicalmachineryinclineparameters}
  \lean{PhyXMiniProblems.ProblemPhyXMini0733.HasPhysicalMachineryInclineParameters}
  \uses{def:physics:phyx-mini-0733:phyxminiproblems-problemphyxmini0733-lengthinmeters, def:physics:phyx-mini-0733:phyxminiproblems-problemphyxmini0733-machineryinclinesetup}
  Positivity and first-quadrant conditions selecting the pictured branch
\end{definition}
\begin{proof}
  This is the declared model definition of Has Physical Machinery Incline Parameters.
\end{proof}
\begin{definition}[Satisfies Inclined Plank Projection Law]
  \label{def:physics:phyx-mini-0733:phyxminiproblems-problemphyxmini0733-satisfiesinclinedplankprojectionlaw}
  \lean{PhyXMiniProblems.ProblemPhyXMini0733.SatisfiesInclinedPlankProjectionLaw}
  \uses{def:physics:phyx-mini-0733:phyxminiproblems-problemphyxmini0733-lengthreadout, def:physics:phyx-mini-0733:phyxminiproblems-problemphyxmini0733-machineryinclinesetup}
  The governing right-triangle projection law. For motion a distance d up a straight plank at angle theta above horizontal, the vertical component is d * sin(theta). The relation is required in every length unit and contains neither the supplied numerical inputs nor the requested answer
\end{definition}
\begin{proof}
  This is the declared model definition of Satisfies Inclined Plank Projection Law.
\end{proof}
\begin{definition}[Answer Choice]
  \label{def:physics:phyx-mini-0733:phyxminiproblems-problemphyxmini0733-answerchoice}
  \lean{PhyXMiniProblems.ProblemPhyXMini0733.AnswerChoice}
  Labels of the four metre-valued answer choices
\end{definition}
\begin{proof}
  This is the declared model definition of Answer Choice.
\end{proof}
\begin{definition}[displayed Vertical Distance Meters]
  \label{def:physics:phyx-mini-0733:phyxminiproblems-problemphyxmini0733-displayedverticaldistancemeters}
  \lean{PhyXMiniProblems.ProblemPhyXMini0733.displayedVerticalDistanceMeters}
  \uses{def:physics:phyx-mini-0733:phyxminiproblems-problemphyxmini0733-answerchoice}
  Metre value printed beside each answer label
\end{definition}
\begin{proof}
  This is the declared model definition of displayed Vertical Distance Meters.
\end{proof}
\begin{definition}[recorded Dataset Answer]
  \label{def:physics:phyx-mini-0733:phyxminiproblems-problemphyxmini0733-recordeddatasetanswer}
  \lean{PhyXMiniProblems.ProblemPhyXMini0733.recordedDatasetAnswer}
  \uses{def:physics:phyx-mini-0733:phyxminiproblems-problemphyxmini0733-answerchoice}
  The answer label recorded in the source dataset
\end{definition}
\begin{proof}
  This is the declared model definition of recorded Dataset Answer.
\end{proof}
\begin{definition}[Rounds To Displayed Hundredth]
  \label{def:physics:phyx-mini-0733:phyxminiproblems-problemphyxmini0733-roundstodisplayedhundredth}
  \lean{PhyXMiniProblems.ProblemPhyXMini0733.RoundsToDisplayedHundredth}
  \uses{def:physics:phyx-mini-0733:phyxminiproblems-problemphyxmini0733-lengthquantity, def:physics:phyx-mini-0733:phyxminiproblems-problemphyxmini0733-lengthinmeters, def:physics:phyx-mini-0733:phyxminiproblems-problemphyxmini0733-answerchoice, def:physics:phyx-mini-0733:phyxminiproblems-problemphyxmini0733-displayedverticaldistancemeters}
  A physical distance agrees with a value displayed to the nearest hundredth of a metre when its metre readout differs by less than half a hundredth
\end{definition}
\begin{proof}
  This is the declared model definition of Rounds To Displayed Hundredth.
\end{proof}
\begin{definition}[Is Unique Rounded Answer]
  \label{def:physics:phyx-mini-0733:phyxminiproblems-problemphyxmini0733-isuniqueroundedanswer}
  \lean{PhyXMiniProblems.ProblemPhyXMini0733.IsUniqueRoundedAnswer}
  \uses{def:physics:phyx-mini-0733:phyxminiproblems-problemphyxmini0733-lengthquantity, def:physics:phyx-mini-0733:phyxminiproblems-problemphyxmini0733-answerchoice, def:physics:phyx-mini-0733:phyxminiproblems-problemphyxmini0733-roundstodisplayedhundredth}
  A choice is the unique answer matching the vertical distance to 0.01 m
\end{definition}
\begin{proof}
  This is the declared model definition of Is Unique Rounded Answer.
\end{proof}
\begin{lemma}[vertical Displacement In Meters eq sine Projection]
  \label{lem:physics:phyx-mini-0733:phyxminiproblems-problemphyxmini0733-verticaldisplacementinmeters-eq-sineprojection}
  \lean{PhyXMiniProblems.ProblemPhyXMini0733.verticalDisplacementInMeters_eq_sineProjection}
  \uses{def:physics:phyx-mini-0733:phyxminiproblems-problemphyxmini0733-lengthinmeters, def:physics:phyx-mini-0733:phyxminiproblems-problemphyxmini0733-anglefromdegrees, def:physics:phyx-mini-0733:phyxminiproblems-problemphyxmini0733-machineryinclinesetup, def:physics:phyx-mini-0733:phyxminiproblems-problemphyxmini0733-matchesmachineryinclinereadouts, def:physics:phyx-mini-0733:phyxminiproblems-problemphyxmini0733-satisfiesinclinedplankprojectionlaw}
  The projection law first gives the exact trigonometric expression for the vertical displacement. This conclusion specializes only the governing law and the two supplied readouts; it is not a premise of the model
\end{lemma}
\begin{proof}
  The conclusion follows from the governing relations and figure readouts named in the statement of vertical Displacement In Meters eq sine Projection.
\end{proof}
% --- Archon declaration topology end ---
% --- Archon physics formalization source end ---
